\section{Grownwallの不等式}

\besha
\begin{theo} (\tf{Grownwallの不等式})
$I=[a,b]$を有界閉区間, $f,g$を$I$上の連続関数で, $g$は非負値であるとする。ある定数$C$に対し,
\[f(t) \leq C+ \int_{a}^{t} f(s)g(s) \, ds\]
がすべての$t\in I$で成り立つと仮定するとき, 次の不等式が成り立つ。
\[ f(t) \leq C\exp{\parena{\int_{a}^{t} g(s)\, ds}} \]
\end{theo}
\ensha
\begin{point}
Liは, これを手で書き殴って丸暗記した。次のような特徴の定理であることを掴もう。
\ben
\item 有界閉区間上。とりあえず連続。
\item 非負値じゃないほうを, 非負値のほうで抑える！
\item \tf{存在性を語る定理ではなく, 一意性を示すときに少し使える定理。}
\een
\end{point}






\begin{egprob} (2019解析学入門演習)\\
$[0,\infty)$上の連続関数$f(x)$により以下の積分表示による微分方程式を考える。
\[y(x) = e^{-2x} + \disp\int_{0}^{x} e^{-2(x-s)} f(s)y(s)\, ds, \quad x\in [0,\infty)\]
$\int_{0}^{\infty} |f(x)|\, dx < \infty$を仮定するとき, この方程式の解$y(x)$は$\disp\lim_{x\to \infty}y(x) = 0$を満たすことを示せ。
\end{egprob}

\subsection{例題}








\clearpage
